\documentclass[11pt]{article}
\usepackage[a4paper,margin=26mm]{geometry}
\usepackage{amsmath,amssymb,mathtools}
\usepackage{unicode-math}
\setmainfont{Latin Modern Roman}
\setmathfont{Latin Modern Math}
\usepackage{xurl,hyperref}
\hypersetup{hidelinks,pdftitle={UnsolvedMath Problem 10300016}}
\newcommand{\sref}[2]{\hyperlink{source:#1:#2}{[S#2]}}
\newcommand{\cataloguescope}[1]{#1}
\begin{document}
% UnsolvedMath ID 10300016; source catalogue code AMR-102-0016
\setcounter{section}{10300015}
\section{Branched surfaces and triangulations — Question 7.1}\label{10300016}
{\small\sffamily UnsolvedMath ID 10300016 · Topology}
\subsection{Definitions and mathematical statement}

\paragraph{Shared notation.}$\mathbb N=\{1,2,\ldots\}$, and $\mathbb Z,\mathbb Q,\mathbb R,\mathbb C$ denote the integers, rationals, reals and complex numbers. Any other conventions are specified locally.


A branched surface $B\subset M$ is a smoothly tangent $2$-complex with the standard local branching models. Its fibered neighborhood $N(B)$ is an interval bundle with the usual identifications over the branch locus. A lamination is carried by $B$ if it lies in $N(B)$ transverse to interval fibers, and fully carried if it meets every fiber. Splitting means cutting this fibered neighborhood along carried surface pieces and taking the induced branched surface.
\paragraph{Statement recorded in the supplied source.}
Characterize embedded branched surfaces $B\subset M$ admitting a nontrivial splitting to a branched surface $B'$ homeomorphic to $B$ as a branched surface. Nontrivial splitting changes the carrying neighborhood by an actual cut, rather than merely applying an ambient isotopy. \sref{10300016}{2}
\subsection{Short English statement}
Characterize branched surfaces that reproduce their own branched topological type after a nontrivial splitting.
\subsection{Sources}
\begin{description}
\item[\hypertarget{source:10300016:1}{S1}] ULAM.AI and the UnsolvedMath Contributors, \emph{UnsolvedMath: A Curated Collection of Open Mathematics Problems}, record 10300016 (AMR-102-0016). CC BY 4.0. \url{https://huggingface.co/datasets/ulamai/UnsolvedMath}.
\item[\hypertarget{source:10300016:2}{S2}] Danny Calegari, \emph{Problems in foliations and laminations of 3-manifolds}, version 0.78 (2002), Question 7.1; Definitions 1.1--1.10 and local definitions and remarks. \url{https://arxiv.org/abs/math/0209081v1}.
\end{description}
\label{end:10300016}
\end{document}
